\subsection{The law of inertia.}

THEOREM 1 (“law of inertia”). --- Suppose that A satisfies the hypotheses at the beginning of this paragraph, and that E is of finite dimension n.~Let Φ be a Hermitian form on E. Then:

\begin{enumerate}
\def\labelenumi{\alph{enumi})}
\item
  There exists a decomposition of E as a direct sum of the subspace \(E^{0}\) orthogonal to E, and of two subspaces \(E^{+}\) and \(E^{-}\) such that the restriction of \(\Phi\) to \(E^{+}\) (resp. \(E^{-}\) ) is positive (resp. negative) and non-degenerate.
\item
  There exists an orthogonal basis \((e_i)_{1\leqslant i\leqslant n}\) of E such that
\end{enumerate}

\[
\Phi (\sum_ {i = 1} ^ {n} \xi_ {i} e _ {i}, \sum_ {i = 1} ^ {n} \eta_ {i} e _ {i}) = \sum_ {i = 1} ^ {s} \xi_ {i} \overline {{\eta}} _ {i} - \sum_ {i = s + 1} ^ {s + t} \xi_ {i} \overline {{\eta}} _ {i}.\tag{2}
\]

\begin{enumerate}
\def\labelenumi{\alph{enumi})}
\setcounter{enumi}{2}
\item
  The dimensions s of \(E^{+}\) and t of \(E^{-}\) are the same for all direct-sum decompositions satisfying the conditions stated in a); the integer s (resp. t) is the maximum of the dimensions of subspaces F of E such that the restriction of \(\Phi\) to F is positive (resp. negative) and non-degenerate.
\item
  The rank of \(\Phi\) is \(s + t\) .
\item
  If \(\Phi\) is non-degenerate, its index is equal to inf(s, t) (§ 4, no. 2, def. 2).
\end{enumerate}

Let, in fact, \((x_{i})\) ( \(i = 1, \ldots, n\) ) be an orthogonal basis of E (§ 6, no. 1, th. 1); recall that one has \(\Phi(x, x) \in K\) for all \(x \in E\) . One may assume that \(\Phi(x_{i}, x_{i}) > 0\) for \(i = 1, \ldots, s\) , \(\Phi(x_{i}, x_{i}) < 0\) for \(i = s + 1, \ldots, s + t\) and \(\Phi(x_{i}, x_{i}) = 0\) for \(i = s + t + 1, \ldots, n\) . This proves \(a\) ), for one takes as \(E^{+}\) (resp. \(E^{-}\) ) the subspace generated by \(x_{1}, \ldots, x_{s}\) (resp. by \(x_{s+1}, \ldots, x_{s+t}\) ), and \(E^{0}\) is then generated by \(x_{s+t+1}, \ldots, x_{n}\) . From this we deduce \(b\) ) by noting that \(\Phi(x_{i}, x_{i})\) is of the form \(\rho\bar{\rho}\) (resp.

\(-\rho\bar{\rho}) \quad (\rho\in\Lambda^{*}) \quad \text{for} \quad i=1,\ldots,s \quad (\text{resp. } i=s+1,\ldots,s+t).\) To prove c), consider a subspace P of E such that the restriction of \(\Phi\) to P is positive and nondegenerate; we then have \(\mathrm{P}\cap(\mathrm{E}^{-}+\mathrm{E}^{0})=\{0\}\) and the sum \(P+E^{-}+E^{0}\) is therefore direct; we conclude that dim \(P\leqslant\dim E^{+}=s\) and this proves c). Assertion d) follows immediately from a).

Finally suppose that \(\Phi\) is nondegenerate, set \(q = \inf (s, t)\) , and let us show that \(q\) is the index of \(\Phi\) . With the notations of \(b\) ), the vectors \(e_i + e_{s+i}\) (resp. \(e_i - e_{s+i}\) ) ( \(i = 1, \ldots, q\) ) generate a totally isotropic subspace F (resp. F'). Since K has characteristic 0, F + F' is generated by \(e_1, \ldots, e_q\) , \(e_{s+1}, \ldots, e_{s+q}\) and is therefore non-isotropic. The restriction of \(\Phi\) to the subspace H = (F + F')⁰ is then positive (or negative) and nondegenerate, and H therefore contains no isotropic vector \(\neq 0\) . Since F⁰ contains F and H, and dim(F + H) = codim F' = codim F, we have F⁰ = F + H. Thus an isotropic vector z orthogonal to F necessarily lies in F, because, in the direct sum F + H, the component of z in H is isotropic. Consequently, F is a maximal totally isotropic subspace, and this proves e).

DEFINITION 2. --- With the notation of Th. 1, the pair \((s,t)\) is called the signature of \(\Phi\) .

